\subsection{Canonical partitions: the precise category and the ambient germ}
\label{sec:A5}

The definitions used here are those of 121/026--028, 030--037 and
049--052.  A manifold has no boundary unless a boundary is explicitly
mentioned.  A product decomposition is an isomorphism in the stated
category, not merely an equality of dimensions or of homology groups.
For a pointed space $(X,x)$, write $\nu_X(x)$ for the largest $d$ for which
there are pointed germs $(N,n)$ and an isomorphism
\[
 (X,x)\simeq(\mathbb R^d\times N,(0,n)).
\]
The axis of such an isomorphism is a germ along which $X$ is
equisingular.  Put $Y_d(X)=\{x:\nu_X(x)=d\}$.
All assertions concerning finite-dimensional polyhedra below use the
PL category and locally finite triangulations.  Statements about a
definable partition specify whether finitely many globally definable
pieces or a locally finite locally definable family is meant.

\begin{saresult}{A5.001}{Refuted}{Proposition A in the unrestricted TOP context}
For locally finite-dimensional topological spaces, Proposition A of
121/027 is false, already for a compact metrizable zero-dimensional
space.
\end{saresult}
\begin{proof}
Let $C\subset[0,1]$ be the middle-thirds Cantor set.  It is totally
disconnected.  If a germ of $C$ had a factor $\mathbb R^d$ with $d>0$,
its product axis would contain a nonconstant path in $C$.  This is
impossible.  Consequently $\nu_C(x)=0$ at every point, so $Y_0(C)=C$.
A zero-dimensional manifold is discrete, whereas $C$ has no isolated
points.  Thus $Y_0(C)$ is not a manifold.  Compactness, metrizability,
and dimension zero do not repair the TOP assertion.  This example
addresses Proposition A; it is not an argument for or against the
separate cancellation assertion in the bracket of Proposition B.
\end{proof}

\begin{saresult}{A5.002}{Refuted}{The recursion of 121/036--037 loses transverse data}
The proposed recursion does not always produce a privileged partition,
even for a proper map of finite compact PL spaces that is a
semialgebraic map.  Therefore it cannot in general produce the
coarsest privileged partition.
\end{saresult}
\begin{proof}
Let $A$ be the simplicial circle with vertices $a_0,\ldots,a_4$ in
cyclic order, and let $B$ be the three-edge circle with vertices
$b_0,b_1,b_2$.  Define a simplicial surjection $g:A\to B$ by
\[
 (g(a_0),g(a_1),g(a_2),g(a_3),g(a_4))
       =(b_0,b_1,b_2,b_1,b_2).
\]
No edge is collapsed.  An interior point of the edge $b_0b_1$ has one
preimage, an interior point of $b_1b_2$ has three, and an interior
point of $b_2b_0$ has one.

Let
\[
 X=M_g=(A\times[0,1])\coprod B\,/
             ((a,0)\sim g(a)),\qquad
 Y=[0,1],\qquad f:X\longrightarrow Y
\]
be the mapping cylinder with its height map, with $f(B)=0$.
It is a proper PL map between finite polyhedra.  Realizing its finite
abstract simplicial complex in a Euclidean space makes the entire
example semialgebraic.  An explicit triangulation has the eight
vertices $a_i,b_j$ and, for $i$ modulo $5$, the two maximal triangles
\[
 [b_{g_i},a_i,a_{i+1}],\qquad
 [b_{g_i},b_{g_{i+1}},a_{i+1}],
 \quad(g_0,g_1,g_2,g_3,g_4)=(0,1,2,1,2).
\]
Here $b_{g_i}$ denotes the target vertex with index $g_i$.

The largest open manifold in $Y$ over which $f$ is locally trivial is
$U=(0,1)$.  Indeed $f^{-1}(U)=A\times(0,1)$, whereas neither endpoint
of $Y$ has an open neighbourhood that is a manifold without boundary.
The space $V=f^{-1}(U)$ is a connected two-manifold.  Its intrinsic
partition therefore has the single leaf $V$, and $\overline V=X$.
The recursive boundary data prescribed in 121/037 consequently are
\[
 Y'=\{0,1\},\qquad X'=B\coprod A,\qquad
 \Phi'=\overline V\cap X'=X'.
\]
Taking finite intersections adds no information.  The recursively
coarsest partition of this marked diagram has the two circle leaves
$B$ and $A$.  The proposed output on $X$ is thus
\[
                  \{V,B,A\}.
\]

It is not equisingular along $B$.  At an interior point of a bottom
edge with $m$ preimages, $X$ is the $m$-page book
$\mathbb R\times c(\{1,\ldots,m\})$.  Its point link is
$S^0*\{1,\ldots,m\}$, the graph $K_{2,m}$.  Therefore
\[
 \operatorname{rank} H_2(X,X\setminus\{p\};\mathbb Q)=m-1.
\]
The rank is $0$ on the interiors of $b_0b_1$ and $b_2b_0$, and $2$
on the interior of $b_1b_2$.  Equisingularity along a connected
manifold leaf makes every local homology rank locally constant on
that leaf, by translating the manifold factor in each product chart.
It would therefore make this rank constant on $B$, a contradiction.

The complete simplicial input and exact rational chain computations
are in \path{code/A5_mapping_cylinder.py} and
\path{logs/A5_mapping_cylinder.json}.  The output is
\[
\begin{array}{c|c|c|c}
\text{bottom edge}&\text{pages}&
 (b_0,b_1)\text{ of its midpoint link}&
 \operatorname{rank}H_2(X,X-p;\mathbb Q)\\\hline
b_0b_1&1&(1,0)&0\\
b_1b_2&3&(1,2)&2\\
b_2b_0&1&(1,0)&0
\end{array}
\]
The complex has $f$-vector $(8,18,10)$ and Betti numbers $(1,1,0)$.
These auxiliary checks certify the supplied finite complex; the
argument above explains why the differing ranks refute the claim.
\end{proof}

The missing datum is the germ of each earlier closed leaf
$\overline F_i$ \emph{in the ambient space} along $X'$, rather than
just its trace $\overline F_i\cap X'$.  In this example the trace is
the entire bottom circle and retains none of the page-number
function.  The example has an elementary privileged refinement:
split $B$ at $b_1,b_2$, use the two complementary open arcs and those
two points, and retain $V$ and the top circle $A$.  The two arcs have
constant page numbers $1$ and $3$; all zero-dimensional leaves are
automatically equisingular, and the remaining charts are product
charts.  The closure condition is immediate from the listed leaves.
This checks that the defect is in the proposed recursion, not in
the existence of a privileged partition for the example.

\begin{saresult}{A5.003}{Proved}{PL suspension cancellation and Proposition D}
For compact PL spaces $K,L$, suspension is cancellative:
$\Sigma K\simeq_{\rm PL}\Sigma L$ implies
$K\simeq_{\rm PL}L$.  If $M$ is a PL manifold of dimension $d$, then
\[
 \nu_{M\times X}(m,x)=d+\nu_X(x),\qquad
 Y_j(M\times X)=M\times Y_{j-d}(X).
\]
The convention is $Y_a(X)=\varnothing$ for $a<0$.
\end{saresult}
\begin{proof}
We include the cancellation argument to isolate exactly the PL input.
The point link $\operatorname{lk}_xX$ is well defined up to PL
homeomorphism: two compatible triangulations have a common
subdivision, and piecewise linear radial projection identifies the
two links.  A pointed PL germ has a factor $\mathbb R^r$ precisely
when its point link is an $r$-fold suspension.  One implication is the
link formula for products; the converse follows by taking the open
cone of the suspension, using
$c(S^{r-1}*K)\simeq_{\rm PL}\mathbb R^r\times cK$ as germs.

First observe that
\begin{equation}\label{eq:A5-desuspend}
 \Sigma^rK\simeq_{\rm PL}\Sigma^sL,\quad r<s
 \quad\Longrightarrow\quad K\text{ is a suspension}.
\end{equation}
Induct on $r$.  The case $r=0$ is the assertion itself.  For $r>0$,
take a point on the suspension sphere $S^{r-1}$ in the first space.
Its link is $\Sigma^{r-1}K$.  The image point in
$S^{s-1}*L$ lies either on $S^{s-1}$, in which case its link is
$\Sigma^{s-1}L$, or off that sphere, in which case its link is
$\Sigma^s\operatorname{lk}_{\ell}L$ for a suitable $\ell\in L$.
In either case the induction hypothesis applies, proving
\eqref{eq:A5-desuspend}.

Desuspend $K$ and $L$ maximally, writing
$K=\Sigma^aK_0$ and $L=\Sigma^bL_0$, where neither $K_0$ nor $L_0$
is a suspension.  This process terminates by dimension; the empty
space is allowed, so that spheres are included.  Applied to
$\Sigma K\simeq\Sigma L$, \eqref{eq:A5-desuspend} forces $a=b$.
Under a homeomorphism
$\Sigma^{a+1}K_0\to\Sigma^{a+1}L_0$, a point of the first suspension
sphere must go to the second suspension sphere: otherwise the link
formula and \eqref{eq:A5-desuspend} would make $K_0$ a suspension.
Taking links cancels one suspension.  Repetition gives $K_0\simeq L_0$
and hence $K\simeq L$.

Finally, if $\operatorname{lk}_xX=\Sigma^qK_0$ with $K_0$ not a
suspension, then
$\operatorname{lk}_{(m,x)}(M\times X)=\Sigma^{d+q}K_0$.
Equation \eqref{eq:A5-desuspend} says that this link admits no
additional suspension.  This proves the two displayed formulas.
The argument is local, so it applies to locally finite polyhedra.
The original checked source is Armstrong's 1966 thesis, paper III,
``Transversality for Polyhedra,'' printed pp.11--17: Definitions
2a--2b and Lemma 1, Lemma 3 with its corollary, and Lemma 4.
Armstrong explicitly credits Lemma 3 and its corollary jointly to
Morton and himself \cite{A5-Armstrong}.
\end{proof}

\begin{saresult}{A5.004}{Proved}{PL Propositions A, B, and C}
Each $Y_i(X)$ is a PL manifold of pure dimension $i$, and $X$ is
PL equisingular along it.  Its transverse pointed PL germ is
immobile and unique up to pointed PL isomorphism, and its type is
constant on each connected component of $Y_i(X)$.  For $x\in Y_i(X)$,
the germ $(Y_i(X),x)$ is the largest manifold germ along which $X$ is
equisingular.  The other such germs are exactly the locally flat
PL submanifold germs of $(Y_i(X),x)$.
\end{saresult}
\begin{proof}
Choose a maximal product germ
$(X,x)\simeq(\mathbb R^i\times N,(0,n))$.  The normal germ has
mobility zero, since any positive Euclidean factor of $(N,n)$ would
increase the displayed factor of $X$.  In a triangulation of $N$
having $n$ as a vertex, every nonvertex point has positive mobility:
the relative interior of its carrier simplex supplies a positive
dimensional manifold factor.  Thus, after shrinking at $n$,
$Y_0(N)=\{n\}$.  By A5.003,
\[
       Y_i(\mathbb R^i\times N)=\mathbb R^i\times\{n\}
\]
in this neighbourhood.  This proves A and the first assertion of B.

For uniqueness, write $N=cL$ as a pointed PL germ.  Two such normal
germs give point-link identifications
$\Sigma^iL\simeq\operatorname{lk}_xX\simeq\Sigma^iL'$.
Repeated cancellation in A5.003 gives $L\simeq L'$, hence
$(cL,*)\simeq(cL',*)$.  In any one product chart the same germ works
at every axis point.  Its now-proved uniqueness makes its type
locally constant, and consequently constant on a connected
component.  No canonical isomorphism, as opposed to a unique
isomorphism class, is asserted.

Let a PL manifold germ $(S,x)$ of dimension $j$ be an equisingular
axis, so that $(X,S)\simeq(\mathbb R^j\times P,
\mathbb R^j\times\{p\})$.  A5.003 implies
$\nu_P(p)=i-j$ and
\[
 Y_i(X)\simeq\mathbb R^j\times Y_{i-j}(P)
\]
near $x$.  The case of A just proved shows that $Y_{i-j}(P)$ is a
manifold of dimension $i-j$ there.  It follows that $S\subset Y_i(X)$
as germs and that this inclusion is locally flat.  Conversely, a
locally flat germ $S\subset Y_i(X)$ admits manifold coordinates
$Y_i(X)\simeq S\times\mathbb R^{i-j}$.  Combining these coordinates
with the product chart of A gives
$X\simeq S\times(\mathbb R^{i-j}\times N)$, which is the required
equisingularity.  This proves C, with ``regularly embedded'' read
as locally flat in the PL category.
\end{proof}

\begin{saresult}{A5.005}{Proved}{The intrinsic partition is privileged and coarsest}
Conjecture 1 of 121/034 holds for locally finite PL spaces and for
definable sets in any fixed o-minimal expansion of the real field,
with isomorphisms in the respective category and with connected
leaves, as in the component definition at 121/031.  The leaves are the
connected components of $Y_i(X)$.  They form a locally finite PL
partition, respectively a finite definable partition.  The partition
is privileged, is coarser than every other privileged partition, and
restricts to opens by taking the connected components of its
intersections with those opens.  In the definable statement, opens
are definable when a finite globally definable partition is required.
\end{saresult}
\begin{proof}
We first give a finite definable argument that also covers sets that
are not locally compact.  Triangulate a bounded realization of $X$
definably as a union of open simplices in a finite complex.  Write
$\mathcal S$ for the finite partition into these open simplices.
Every point of a fixed $S\in\mathcal S$ has the same pointed
definable local germ, by the elementary product chart at the relative
interior of a simplex.  This remains true when some faces of the
complex have been omitted from $X$.

Declare two points equivalent if their pointed neighbourhood germs
in $X$ are definably isomorphic.  Each equivalence class is a union
of members of the finite set $\mathcal S$.  In particular no
quantification over definable functions is being used to claim
definability: there are only finitely many possible classes, each
already a finite union of definable pieces.  Let $r(x)$ be the
largest dimension of a member of $\mathcal S$ in the class of $x$.
If $x$ has a definable $d$-dimensional product axis, all points of
that axis have the same local germ as $x$.  The axis lies in the
finite union of simplices in this class.  Definable dimension gives
$d\leq r(x)$.  Conversely, a point in a simplex of dimension $r(x)$
has an $r(x)$-dimensional product factor; transporting its pointed
chart to $x$ shows $\nu_X(x)\geq r(x)$.  Hence
\begin{equation}\label{eq:A5-mobility-types}
                         \nu_X(x)=r(x).
\end{equation}

For $\nu_X(x)=i$, choose an equivalent point $z$ in an $i$-simplex
of $\mathcal S$.  In the product chart about $z$, its carrier
simplex is the axis $B^i\times\{n\}$, with $B^i$ an open ball.
Every off-axis point is in a carrier simplex of dimension greater
than $i$, and hence has mobility greater than $i$.  Every axis
point has mobility $i$.  Transporting the chart to $x$ proves
\[
              (X,Y_i(X))\simeq(B^i\times N,B^i\times\{n\})
\]
locally.  Thus $Y_i(X)$ is a definable $i$-manifold and has finitely
many connected components.  Notice that this proof did not assert
PL cancellation for arbitrary noncompact semilinear sets and did not
use a compact version of the Hauptvermutung outside its scope.

The PL case follows either from A5.004 or from the same finite local
argument.  Mobility is constant on each open simplex of every
locally finite triangulation.  The sets $\{\nu\leq i\}$ are closed:
the complementary set is open because a product chart with a factor
$\mathbb R^{i+1}$ gives that factor at every point of the chart.
Consequently these closed sets are subcomplexes.  There are locally
only finitely many connected components of the differences, so the
proposed PL partition is locally finite.

It remains to verify the full diagram condition in the definition
of privileged, not merely equisingularity of $X$.  In a product
chart $B^i\times N$ along one intrinsic leaf, translation in $B^i$
preserves the pointed local germ of every point.  Therefore it
preserves mobility.  If $F$ is any global intrinsic leaf, a path in
$B^i\times\{n'\}$ stays in one $Y_j(X)$ and hence in one connected
component.  Thus
\[
 F\cap(B^i\times N)=B^i\times E_F,
 \qquad
 \overline F\cap(B^i\times N)=B^i\times\overline{E_F}
\]
for a subset $E_F\subset N$, with the closures in this formula
relative to the chart.  This holds simultaneously for all the
locally finitely many incident leaves.  It proves the required
product trivialization of the entire inclusion diagram of leaf
closures.  It also shows that the intersection of $\overline F$
with any other connected intrinsic leaf is open as well as closed
in that leaf; hence it is empty or the whole leaf.  All leaves are
locally closed, since each $Y_i$ is a difference of closed sets and
its connected components are open and closed in it.  These facts
are exactly the three conditions of 121/032 together with the
manifold requirement of 121/034.

Finally let $S$ be a leaf of any privileged partition.  Its product
charts imply that the local germ class, and therefore mobility, is
locally constant on $S$.  Since $S$ is connected, it lies in one
$Y_i$ and in one connected component of that $Y_i$.  This proves
coarseness.  All constructions use local germs, so an open
restriction gives the same mobility sets; splitting their
intersections into connected components gives precisely the stated
restriction rule.
\end{proof}

The definable proof uses the usual triangulation theorem, rather
than presuming that a definable homeomorphism is itself PL.  For
compact PL pairs, the separate bridge is Shiota's o-minimal
Hauptvermutung, Theorem 2.1, p.5, of \cite{A5-Shiota}: definably
homeomorphic compact polyhedral pairs are PL homeomorphic.  That
theorem and the elementary local argument agree on locally compact
definable germs.  The noncompact qualifications in the same paper,
pp.22--24, explain why one must not state global semilinear
cancellation without hypotheses.

\begin{saresult}{A5.006}{Refuted}{The algebraic context of 121/049 with classical local finiteness}
In the category of complex algebraic varieties with globally
locally closed algebraic leaves, the assertion of 121/049 is false
when ``locally finite'' refers to the classical topology.  It fails
at the existence stage.  A locally finite closed family likewise
need not admit a finite globally definable adapted partition.
\end{saresult}
\begin{proof}
Take $X=\mathbb A^1_{\mathbb C}$, with its classical topology, and
$T_n=\{n\}$ for $n\in\mathbb Z$.  This is a locally finite family
of closed algebraic subspaces.  If a partition is adapted to each
$T_n$, the leaf containing $n$ must be precisely $\{n\}$.
Every positive-dimensional locally closed algebraic subset of
$\mathbb A^1_{\mathbb C}$ is cofinite: its Zariski closure is the
whole affine line, and it is open in that closure.  Such a subset
contains infinitely many integers, so it cannot be a leaf of an
adapted partition.  All algebraic leaves must consequently be
zero-dimensional; if leaves are connected, they are singletons.
The singleton partition of $\mathbb C$ is not locally finite.
This proves nonexistence even before the strict condition is imposed.

For finite globally definable partitions, $X=\mathbb R$ with the
same family $T_n=\{n\}$ is enough: every $T_n$ must be a separate
leaf, requiring infinitely many leaves.  Both counterexamples use
an infinite locally finite family, rather than a finite input whose
certificate is omitted.
\end{proof}

If ``locally finite'' in the algebraic context instead means
Zariski locally finite, quasi-compactness forces only finitely many
nonempty $T_\alpha$.  The counterexample does not apply to that
different statement.  If leaves are permitted to be locally
analytic or locally semialgebraic, it also does not apply: those
are larger categories.  These distinctions are necessary to give
121/049 a precise truth value.

\begin{saresult}{A5.007}{Proved}{Strict admissibility and the forced first open set}
Read ``strictly admissible'' as 121/051: the leaves are connected
and locally connected, the partition is locally finite and satisfies
the frontier condition, and for every open $W\subset X$, the
connected components of all $F\cap W$ again give a locally finite
frontier partition of $W$.  Every locally finite PL or locally definable conically
equisingular manifold partition is strict.  A finite family of
closed PL subspaces, respectively closed definable subspaces,
therefore has a strict adapted partition.  For any strict adapted
partition the set
\[
 U_0=X\setminus\bigcup_\alpha\partial_XT_\alpha
\]
is saturated and open, and its connected components cannot be
crossed by a leaf.  The same statement holds locally for a locally
finite family.  The draft expression
$\bigcup_\alpha\operatorname{int}_X(T_\alpha)$ in 121/049 is not
the required first open set.
\end{saresult}
\begin{proof}
For a conical partition, fix $x\in W$ and take a small product-cone
chart contained in $W$.  It meets only finitely many leaves, each
in finitely many connected pieces.  Any component of $F\cap W$
meeting a still smaller such chart contains one of those pieces;
this proves the required local finiteness.  If a component $C$ of
$F\cap W$ has a point of another local leaf in its closure, the
product chart along that leaf shows that the closure contains an
open neighbourhood of the point in that leaf: the manifold-factor
paths stay in $C$.  Thus the intersection of $\overline C^{\,W}$
with any component of any other leaf in $W$ is both open and
closed.  Connectedness gives the frontier condition.  Triangulating
the finite closed family simultaneously supplies such a partition
by relative interiors of simplices.  Finite open-simplex models
give the identical argument in the definable setting.

For the saturation claim, in any locally finite frontier partition,
the closure of a saturated subset is saturated: it is locally a
finite union of closures of leaves.  For closed saturated $T$, its
interior equals $X\setminus\overline{X\setminus T}$, so is
saturated as well.  Hence $\partial_XT=T\setminus\operatorname{int}T$
is saturated.  Locally finite unions preserve closedness and
saturation, proving the claim for $U_0$.  Each connected leaf in
$U_0$ lies in one component of $U_0$.  On $U_0$ all membership
conditions in the $T_\alpha$ are locally constant.

Finally, for $X=\mathbb R$ and the single marked subset
$T=\{0\}$, the draft union of interiors is empty, whereas the
first available open region is $\mathbb R\setminus\{0\}$.
The formulas differ already for this single marked subset.
\end{proof}

\begin{saresult}{A5.008}{Proved under stated hypotheses}{The maximal triviality open: the part supplied by Hardt and the triangulable case}
For every continuous definable map $f:X\to Y$, there is a dense
definable open subset $U_0\subset Y$ which is a manifold, possibly
with components of different dimensions, and over which $f$ is
definably locally trivial.  Consequently the maximal such open
$U$ exists and has nowhere dense complement.  If $\dim Y\leq1$,
then $U$ is definable.  More generally, $U$ is definable if $f$
admits a finite definable triangulation as a simplicial map; this
includes proper finite PL maps.  Properness is not needed for
the first assertion.
\end{saresult}
\begin{proof}
Use Hardt's theorem for definable families in the form proved in
Coste, Theorem 5.22, pp.62--63 of \cite{A5-Coste}.  Applied to
the graph of $f$, it gives a finite definable partition of $Y$ over each member of
which $f$ is trivial, and refine it into a finite frontier partition
by definable manifolds.  The union $U_0$ of its open members is
open and dense: the frontier of every member has smaller dimension,
and locally maximal members cover a dense open set.  Over each
such member the Hardt trivialization restricts to product charts.
The union of all open manifold neighbourhoods with such charts is
the required maximal open $U$.  It contains $U_0$, so
$Y\setminus U\subset Y\setminus U_0$ is nowhere dense.  For a
definable curve, $Y\setminus U_0$ is finite.  Any subset of a finite
set is definable with parameters, so $U$ is definable in that case.

Suppose next that $f:|K|\to|L|$ is simplicial with $K,L$ finite.
For two points in the relative interior of one target simplex
$\sigma$, the pointed local map diagrams are isomorphic.  To
check the whole preimage, use barycentric coordinates.  For each
source vertex mapping to a vertex $v$ of $\sigma$, multiply its
barycentric coordinate by the same positive factor used for the
$v$-coordinate in a small target reparametrization; leave other
vertex groups unchanged and normalize the total sum.  The target
operation and its source lift commute with $f$, have inverse
operations of the same kind, and agree on faces.  They therefore
define simultaneous semialgebraic homeomorphisms on a neighbourhood
of the full inverse image of the chosen target neighbourhood.
This shows that manifoldhood of $Y$ and local triviality of $f$
are constant along each open target simplex.  Their common locus
$U$ is a union of finitely many such simplices and is definable.
Transport this conclusion through the assumed definable
triangulation.
\end{proof}

Hardt's theorem alone does not identify the maximal open $U$ with
the union of the open members of an arbitrarily chosen Hardt
partition.  That union can omit regular points on artificial
partition boundaries.  Thus the general higher-dimensional
definability assertion is kept as a separate gate below; it is
not silently deduced from Hardt.  Regardless of that gate, A5.002
disproves the asserted recursive conclusion for a proper map
whose maximal open is explicitly definable.

\begin{saresult}{A5.009}{Refuted}{No coarsest strict partition, already for a compact finite PL pair}
The claim of 121/049 is false in the PL and semialgebraic contexts,
even for one closed marked subspace of a compact finite polyhedron.
Here ``strict'' has exactly the connectedness, local connectedness,
and universal open-restriction meaning of 121/051, conditions (C),
not the additional manifold or equisingularity conditions listed
separately on that page.
\end{saresult}
\begin{proof}
Put
\[
 X=\{(x,y,z)\in[-2,2]^2\times[0,2]:
            x\leq0\ \text{or}\ y\geq z\ \text{or}\ y\leq-z\},
 \qquad T=X\cap\{z=0\}.
\]
These are compact finite polyhedra and $T=[-2,2]^2\times\{0\}$
is closed.  Define locally closed PL subspaces
\[
\begin{aligned}
 A&=X\setminus T,&
 T_+&=T\cap\{y>0\},&T_-&=T\cap\{y<0\},\\
 S&=T\cap\{y=0\},&
 R&=S\cap\{x\geq0\},&B&=T\setminus R.
\end{aligned}
\]
Consider the two partitions
\[
                \mathcal P=\{A,T_+,T_-,S\},\qquad
                \mathcal Q=\{A,B,R\}.                 \tag{A5.9.1}
\]
Every displayed leaf is connected and locally connected.  In
particular $B$ is a rectangle with a slit from its centre to the
right edge, and $R$ is the closed slit, including its endpoint at
$(0,0,0)$.  Manifoldhood without boundary is not part of strictness:
121/051 explicitly lists that as the separate condition (b), whereas
its condition (C) imposes connectedness, local connectedness, local
finiteness of components, and their frontier condition.  Ordinary
admissibility is defined at 121/052.  The exact transcription
locators are lines 93--112 and 141--148 of
\texttt{MPL\_121\_049-073+153+155\_transcription.tex}.

The global closures already show ordinary admissibility:
\[
 \overline A=X,\quad \overline{T_\pm}=T_\pm\cup S,
 \quad\overline S=S,\quad\overline B=T,\quad\overline R=R.
                                                               \tag{A5.9.2}
\]
We verify strictness for \emph{every} open set, including
nondefinable ones.  We use the following elementary local-to-global
criterion.  Suppose a finite tame partition has, at every point, an
arbitrarily small neighbourhood basis in which the connected
components of its leaves form a finite frontier partition.  Then
it is strict.  Indeed, for any open $W$ and $p\in W$, choose one
such neighbourhood $V\subset W$.  Every component of a leaf in $W$
meeting $V$ contains a component of that leaf in $V$, proving local
finiteness near $p$.  If $C$ is a component in $W$ and
$p\in\overline C^{\,W}\cap D$ for another component $D$, the
finitely many local components in $V$ show that one component
$C_0\subset C\cap V$ has $p$ in its closure.  The local frontier
condition puts the component of $D\cap V$ through $p$ in
$\overline{C_0}^{\,V}$.  Local connectedness makes this a
neighbourhood of $p$ in $D$.  Hence $\overline C^{\,W}\cap D$ is
both open and closed in connected $D$, and is all of $D$ if nonempty.
Components are locally closed because leaves are locally closed
and locally connected.  This proves the criterion.

Here the neighbourhoods can be chosen as small coordinate boxes,
intersected with $X$.  At a point of $S$ with $x>0$, the two local
components of $A$ are the wings $y\geq z$, $z>0$, and $y\leq-z$, $z>0$.  Their closures meet $T$ in the upper and
lower closed half-boxes, respectively.  Each contains all of the
local $S$, and each contains precisely its corresponding local
$T_+$ or $T_-$.  At a point of $S$ with $x<0$, $A$ is a connected
upper half-box and its closure contains all local $T$.  At $x=0$,
the two wings are joined by the part $x<0,z>0$ in every small box;
again $A$ is connected there and its closure contains all local
$T$.  Finally, at a point of $T_+$ or $T_-$, take the box so small
that $|y|>z$ throughout its part with $z>0$.  There $A$ is
connected and closes onto the whole bottom face.  The parts on the
outer boundary of $[-2,2]^2\times[0,2]$ are the corresponding
half-boxes and satisfy the identical statements.  These descriptions,
together with $\overline{T_\pm}=T_\pm\cup S$, give a finite
frontier partition in every such box for $\mathcal P$.

For $\mathcal Q$, at a point of $R$ with $x>0$, the local $B$ has
two components, the upper and lower half-boxes, and the two $A$
wings close onto those components and the entire local $R$.
At the endpoint $x=0$, both the local slit complement $B$ and the
local $A$ are connected.  Their closures contain the whole local
$R$.  At a point of $B$, choose a box missing $R$.  If $y\ne0$,
then $|y|>z$ can again be arranged; if $y=0$, then $x<0$ and the
whole upper half-box is present.  In either case the local $A$ is
connected and closes onto the entire local $B$.  Points of $A$
have neighbourhoods contained in $A$.  Thus the criterion proves
strictness of both partitions in (A5.9.1).

Any common coarsening of $\mathcal P$ and $\mathcal Q$ adapted
to $T$ must be exactly $\{A,T\}$.  Both $T_+$ and $T_-$ lie in
$B$, and $S$ meets both $B$ and $R$, so all of $T$ must merge.
Adaptation to $T$ forbids merging any part of it with $A$.
But $\{A,T\}$ is not strict.  In a small box about $(1,0,0)$,
$T$ is connected while $A$ has two components.  The positive wing
has closure meeting $T$ in its upper closed half, a nonempty proper
subset of that connected local leaf.  The required frontier
condition fails.  A coarsest strict adapted partition would be a
common coarsening of (A5.9.1), which is impossible.

For an independently reproducible finite certificate, triangulate
the $(y,z)$ rectangle using vertices
\[
 (-2,0),(0,0),(2,0),(-2,2),(0,2),(2,2)
\]
and triangles $013,134,145,125$.  Over $x\in[-2,0]$ take all four
triangular prisms; over $x\in[0,2]$ take only $013$ and $125$.
The common staircase subdivision gives 18 tetrahedra, with
$f$-vector $(17,50,52,18)$ and exact total volume $24$.
The program \path{code/A5_strict_no_coarsest.py} constructs
all faces and the leaf labels, computes connected components by
the face-poset comparability graph, and checks their relative
closures in every simplex open star.  Its saved output is
\begin{verbatim}
P: 137 simplex-star checks; failures = 0
Q: 137 simplex-star checks; failures = 0
join: 137 simplex-star checks; failures = 2
positive-axis carrier = [(0, 0, 0), (2, 0, 0)]
PASS: P and Q are strict; their only T-adapted
      common coarsening is not strict.
\end{verbatim}
The full coordinates, tetrahedra, all 137 simplex labels, and a
failed frontier inclusion are in
\path{logs/A5_strict_no_coarsest.json}; the complete concise
run log is \path{logs/A5_strict_no_coarsest.txt}.
An open star at a carrier simplex has arbitrarily small
stratum-preserving PL copies about every point of that carrier,
so this also supplies a finite certificate of the local criterion.
The preceding box proof establishes the claim independently of
the program.
\end{proof}

\begin{saresult}{A5.010}{Refuted}{The strict-partition claim also fails for pure complex algebraic and analytic spaces}
There is a reduced affine complex algebraic space $X$ of pure
complex dimension three, with one closed algebraic subspace
$T\subset X$, having two strict adapted partitions and no strict
common adapted coarsening.  The failure occurs on a Zariski open
subset.  Consequently the assertion of 121/049 fails in both
$\mathrm{Alg}_{\mathbb C}$ and $\mathrm{An}_{\mathbb C}$, as well
as in the PL and semialgebraic contexts of A5.009.
\end{saresult}
\begin{proof}
Use coordinates $(u,v,w,t,s)$ on $\mathbb C^5$ and set
\[
\begin{aligned}
 Y&=\{t=s=0\},&
 E&=\{v=0,\ t=uw\},& X&=Y\cup E,\\
 T&=\{t=s=w=0\},&
 S_1&=\{t=s=w=v=0\},&
 S_2&=\{t=s=w=u=0\},\\
 S&=S_1\cup S_2,& A&=X\setminus T.
\end{aligned}                                                    \tag{A5.10.1}
\]
We give $X$ its reduced structure.  The maps
\[
\begin{aligned}
 \phi_Y:\mathbb C^3&\longrightarrow Y,&
 (p,q,r)&\longmapsto(p,q,r,0,0),\\
 \phi_E:\mathbb C^3&\longrightarrow E,&
 (p,q,r)&\longmapsto(p,0,q,pq,r)
\end{aligned}
\]
are polynomial isomorphisms, with inverses the restrictions of the
coordinate projections $(u,v,w,t,s)\mapsto(u,v,w)$ and
$(u,v,w,t,s)\mapsto(u,w,s)$, respectively.  Thus $Y$ and $E$ are
smooth irreducible threefolds, neither contained in the other, and
$X$ has pure dimension three.  Also $T\simeq\mathbb C^2$ with
coordinates $(u,v)$; $S_1$ and $S_2$ are its two coordinate axes.

The intersection responsible for the example is
\[
 Y\cap E=\{t=s=v=0,\ uw=0\}=S_1\cup C,
 \qquad C=\{u=v=t=s=0\},                         \tag{A5.10.2}
\]
where $C$ is the $w$-axis.  Hence
\[
 E\cap T=S_1,\qquad
 (Y\setminus T)\cap(E\setminus T)=C\setminus\{0\}.
                                                               \tag{A5.10.3}
\]
The two pieces of $A$ are connected:
\[
 Y\setminus T\simeq\mathbb C^2\times\mathbb C^*,
 \qquad E\setminus T\simeq
       \mathbb C\times(\mathbb C^2\setminus\{0\}).
\]
Their intersection in (A5.10.3) is nonempty, so $A$ is connected.
All spaces in these formulas are locally connected.

Write $B_0=T\setminus S$ and $B_1=T\setminus S_1$, and define
\[
 \mathcal P=\{A,B_0,S\},\qquad
 \mathcal Q=\{A,B_1,S_1\}.                         \tag{A5.10.4}
\]
These are finite partitions into connected, locally connected,
locally closed algebraic subspaces.  Indeed
$B_0\simeq(\mathbb C^*)^2$, $B_1\simeq\mathbb C\times\mathbb C^*$,
and the two lines in $S$ meet at the origin.  Their closures satisfy
\[
 \overline A=X,\quad\overline{B_0}=\overline{B_1}=T,
 \quad\overline S=S,\quad\overline{S_1}=S_1,
                                                               \tag{A5.10.5}
\]
so both partitions are admissible and adapted to $T$.

We check strictness in the classical topology, which also implies
the required assertion on Zariski open subsets.  The local-to-global
criterion proved in A5.009 applies to these semialgebraic spaces.
At the origin take ambient polydiscs with a common radius
$0<\varepsilon<1$.  In $Y$ the complement of $T$ is
$D_u\times D_v\times D_w^*$, which is connected.  In $E$ it is
\[
 D_u\times\bigl((D_w\times D_s)\setminus\{(0,0)\}\bigr).
\]
The extra ambient inequality $|t|=|uw|<\varepsilon$ is automatic
because $|u|,|w|<\varepsilon<1$.  This piece is connected, and
it meets the first piece in the punctured $w$-disc of $C$.
Thus $A$ has a single local connected component at the origin,
and its closure contains all local $T$.  The local $B_0$ is
$D_u^*\times D_v^*$, the local $B_1$ is $D_u\times D_v^*$,
and the local $S$ is the union of the two coordinate discs.
Each is connected and has the closure stated in (A5.10.5),
relative to the polydisc.  Hence both local partitions have the
frontier property there.

At $p\in S_1\setminus\{0\}$, keep $u$ away from zero and use
instead the holomorphic coordinate change
\[
 \psi:\{u\ne0\}\longrightarrow\{u\ne0\},\qquad
 (u,v,w,t,s)\longmapsto(u,v,uw,t,s).
\]
Its inverse divides the third coordinate by $u$.  In the resulting
coordinates $w'=uw$, the equations are $Y:\ t=s=0$ and
$E:\ v=0,t=w'$.  Small coordinate polydiscs show that the local
$A$ has exactly two connected components:
\[
 A_Y=Y\setminus T,\qquad A_E=E\setminus T,
\]
restricted to that neighbourhood.  They are disjoint because
$Y\cap E$ there is $S_1\subset T$.  The first is a product with a
punctured complex disc and closes onto all local $T$; the second
is a product with a punctured complex two-disc and closes onto
precisely the local $S_1$.  For $\mathcal P$, the local $S$ equals
$S_1$, and for $\mathcal Q$ it is already the designated leaf.
The local $B_0=B_1$ is connected and dense in $T$.  Thus both
partitions again satisfy the local frontier condition.

At $p\in S_2\setminus\{0\}$, $v\ne0$ and a small neighbourhood
misses $E$.  We have only the smooth component $Y$, its divisor
$T$, and, for $\mathcal P$, the coordinate axis $S_2$ in $T$.
Products of discs and punctured discs verify all the stated local
connectedness and frontier conditions.  The same, with fewer
leaves, applies at points of $B_0$.  Points of $A$ admit connected
neighbourhoods contained in $A$, since $A$ is open and locally
connected.  These cases exhaust $X$.  Applying the criterion proves
that both partitions in (A5.10.4) remain admissible after restricting
to \emph{any} classical open set and splitting into connected
components.  This establishes full strictness, not merely a check
on a selected finite collection of opens.

Any common coarsening of the two partitions adapted to $T$ must
be $\{A,T\}$.  The leaf $B_1$ contains $B_0$ and meets $S$ in
$S_2\setminus\{0\}$, so a common coarsening merges $B_0$ and $S$;
these are all of $T$.  Adaptation forbids merging $T$ with $A$.
But take the Zariski open subset
\[
                         U=X\cap\{u\ne0\}.
\]
The two connected components of $A\cap U$ are
\[
 (Y\setminus T)\cap U\simeq\mathbb C^*\times\mathbb C\times\mathbb C^*,
 \qquad
 (E\setminus T)\cap U\simeq
       \mathbb C^*\times(\mathbb C^2\setminus\{0\}).
\]
They are disjoint by (A5.10.2).  The closure of the second meets
$T\cap U\simeq\mathbb C^*\times\mathbb C$ in
$S_1\cap U\simeq\mathbb C^*$, a nonempty proper subset of that
connected local leaf.  Thus $\{A,T\}$ fails strictness on the
Zariski open $U$.  No coarsest strict adapted partition exists.

The program \path{code/A5_complex_strict_no_coarsest.py}
verifies the finite polynomial data by exact Gr\"obner elimination
over $\mathbb Q$, so the identities persist over $\mathbb C$.
It checks all intersections in (A5.10.2)--(A5.10.3), both
parametrizations with their inverse projections, the coordinate-axis
union, and the localization that separates the two branches.  In
particular its output includes
\begin{verbatim}
X = union of two smooth affine 3-spaces;
    pure complex dimension 3.
I(Y intersection E) = (t, s, v, u*w)
I(E intersection T) = (t, s, v, w)
I(S1 union S2) = (t, s, u*v, w)
I((Y intersection E) localized at w) = (t, s, v, u)
I((Y intersection E) localized at u*w) = (1,)
Exact ideal and parametrization checks: 12 passed.
Local connected-piece frontier patterns: 6 passed.
PASS: algebraic identities and incidence certificate
      verified exactly.
\end{verbatim}
Here ``localized'' denotes the contraction of the localized ideal,
computed by saturation.  The complete ideals, rational points
witnessing the forced common coarsening, frontier tables, and log
are saved in \path{logs/A5_complex_strict_no_coarsest.json}
and \path{logs/A5_complex_strict_no_coarsest.txt}.
The code checks the algebra and incidence tables; connectedness of
the local pieces and the universal-open assertion are proved above,
not inferred from sample points.
\end{proof}

\subsection{The TOP uniqueness bracket: an axis-preserving counterexample}
\begin{saresult}{A5.011}{Refuted}{A maximal TOP axis does not determine its transverse germ}
The uniqueness bracket of Proposition B, 121/027, fails even when
$X$ itself is a locally finite polyhedron, the maximal mobility
stratum is a line, and both competing transverse spaces are locally
compact, finite-dimensional ENRs with exactly one nonmanifold point.
More precisely, let
\[
 L=S^5/C_{23},\qquad
 \zeta\cdot(z_1,z_2,z_3)=(\zeta z_1,\zeta z_2,\zeta z_3),
 \quad \zeta=e^{2\pi i/23}.
\]
There is an open smooth six-manifold $W$ with two ends and an end
compactification $C=W\cup\{v\}$ at its positive end such that
\[
 F:(\mathbb R\times C,\mathbb R\times\{v\})
       \xrightarrow{\ \cong\ }
   (\mathbb R\times cL,\mathbb R\times\{o\}),
 \qquad F(x,v)=(x,o),                                      \tag{A5.11.1}
\]
but $(C,v)$ and $(cL,o)$ are not homeomorphic as germs.
Their vertices are immobile, and every point of the distinguished
axis in (A5.11.1) has mobility exactly one.
\end{saresult}
\begin{proof}
We give the projective module, the finite handle input, the end
obstruction, and the extension across the axis.  Throughout, the
involution on a group ring sends each group element to its inverse;
on $K_0$ it also takes the module dual.  Confusing this involution
with complex conjugation alone changes the computation.

\paragraph{1. An explicit projective module over $\mathbb Z[C_{23}]$.}
Put
\[
 R=\mathbb Z[g]/(g^{23}-1),\quad
 N=1+g+\cdots+g^{22},\quad
 r=1+\sum_{a\in Q}g^a.
\]
Here
\[
 Q=\{1,2,3,4,6,8,9,12,13,16,18\}.
\]
Convolution in $R$ gives
\[
 r^2-r+6=6N,\qquad \overline r=N+1-r.                     \tag{A5.11.2}
\]
Let $P=(2,r)\subset R$.  Multiplication by
$a=(1-r)/2\in R\otimes\mathbb Q$ carries $P$ into $R$,
since $2a=1-r$ and $ar=3(1-N)$.  The two maps
\[
 u:R^2\longrightarrow P,\quad u(b,c)=2b+rc,
 \qquad
 j:P\longrightarrow R^2,\quad j(p)=(ap,p)
\]
satisfy $uj=\operatorname{id}_P$.  Consequently $P$ is projective,
and $ju$ is the explicit idempotent
\[
 E=\begin{pmatrix}1-r&3(1-N)\\2&r\end{pmatrix},\qquad E^2=E.
                                                               \tag{A5.11.3}
\]
Its trace is one and determinant zero, so it has rank one at every
prime.  Set $x=[P]-[R]\in\widetilde K_0(R)$.

\paragraph{2. A proof that $x+x^*$ is nonzero.}
Let $O=\mathbb Z[\zeta]$, where $\zeta$ is a primitive $23$rd root.
The ring map $R\to O$, $g\mapsto\zeta$, carries $P$ to
\[
 J=(2,\omega),\qquad
 \omega=1+\sum_{a\in Q}\zeta^a,
 \qquad \omega^2-\omega+6=0.
\]
More precisely, multiplication gives an isomorphism
$P\otimes_R O\xrightarrow{\sim}J$: its image is $J$, its source
is projective of rank one over the domain $O$, and it is an
isomorphism over the fraction field.  Its kernel is therefore
torsion in a torsion-free module and is zero.
The polynomial $\Phi_{23}(T)$ is irreducible because
$\Phi_{23}(T+1)$ is Eisenstein at $23$.  The displayed expression
for $\omega$ is not rational, so
$K=\mathbb Q(\omega)=\mathbb Q(\sqrt{-23})$ and
$[\mathbb Q(\zeta):K]=11$.  The ring of integers of $K$ is
$O_K=\mathbb Z[\omega]$, its discriminant being the fundamental
discriminant $-23$.

The ideal $\mathfrak p=(2,\omega)\subset O_K$ has norm two.
It is not principal: an element $a+b\omega$ has norm
\[
 a^2+ab+6b^2,\qquad
 4\operatorname N(a+b\omega)=(2a+b)^2+23b^2,
\]
and this cannot equal two.  Indeed $b\ne0$ makes the right side
at least $23>8$, and $b=0$ would give $a^2=2$.
On the other hand,
\[
 \mathfrak p^3=(2-\omega),\qquad
 \mathfrak p\,\overline{\mathfrak p}=(2).                 \tag{A5.11.4}
\]
For the first equality, $2-\omega=8+4\omega+\omega^3$
belongs to $\mathfrak p^3$, and both ideals have norm eight.
For the second, both sides have norm four and the containment
follows by multiplying the two generators of each ideal.
Thus $[\mathfrak p]$ has order exactly three in
$\operatorname{Pic}(O_K)$.

The extended ideal $J=\mathfrak p O$ has order exactly three in
$\operatorname{Pic}(O)$.  To see that it does not become principal,
regard $O$ as a projective $O_K$-module of rank eleven.  If
$\mathfrak p O\simeq O$ as $O$-modules, taking determinants over
$O_K$ would give
\[
 \mathfrak p^{11}\otimes_{O_K}\det_{O_K}O
       \simeq\det_{O_K}O,
\]
contradicting the order three of $[\mathfrak p]$.  This argument
uses only that a finite torsion-free module over the Dedekind ring
$O_K$ is projective; it does not assume a numerical cyclotomic
class-number table.

Base change sends $P^*=\operatorname{Hom}_R(P,R)$, with the
involution-twisted module structure, to
$\overline J^{-1}$.  From $J\overline J=(2)O$ we get
$[\overline J^{-1}]=[J]$.  Taking determinants after base change
therefore gives
\[
 \widetilde K_0(R)\longrightarrow\operatorname{Pic}(O),
 \qquad x+x^*\longmapsto2[J]\ne0.                        \tag{A5.11.5}
\]
In particular $x+x^*\ne0$.

\paragraph{3. The finite smooth construction and the sign.}
Write $z$ for the generator of a second infinite cyclic group.
The matrix
\[
 A(z)=1-E+zE\in\operatorname{GL}_2(R[z,z^{-1}]),\qquad
 A(z)^{-1}=1-E+z^{-1}E                                  \tag{A5.11.6}
\]
represents the Bass summand $b(x)$ in
$\operatorname{Wh}(C_{23}\times\mathbb Z)$.
The rank-one free correction contributes the class of $z$, which
is zero in the Whitehead group.  The Bass projection
$\delta:\operatorname{Wh}(C_{23}\times\mathbb Z)
\to\widetilde K_0(R)$ sends $b(x)$ to $x$.
These assertions can also be read directly from the decomposition
$R^2=P\oplus\ker u$: $A$ is $z$ on $P$ and the identity on the
complement.  The split projection and its geometric infinite-cover
interpretation are stated in \cite[printed pp.124--125,
Example 11.5(ii)]{A5-HughesRanicki1996}.

Let $M=L\times S^1$, a closed oriented six-manifold.
Milnor's realization theorem supplies an oriented seven-dimensional
smooth $h$-cobordism $(Z;M,N')$ with
\[
 \tau(Z,M)=b(x).                                         \tag{A5.11.7}
\]
One may use two index-two handles and two index-three handles with
incidence matrix $A(z)$; the construction and the embedded-sphere
argument are \cite[\S11, Existence Theorem 11.1 and Lemma 11.2,
pp.398--399]{A5-Milnor1966}.  Here and below $N'$ denotes a closed
six-manifold, not the group-ring element $N$ of (A5.11.2).
All fundamental groups are identified through the inclusions into
$Z$ and the retraction $Z\to M$.

For a seven-dimensional cobordism, Milnor duality has sign
$(-1)^{7-1}=+1$, so
\[
 \tau(Z,N')=\overline{b(x)}=-b(x^*).
\]
The notation $\overline{b(x)}=-b(x^*)$ means the involution on
$\operatorname{Wh}(C_{23}\times\mathbb Z)$, which inverts $z$ as
well as $g$.  The induced homotopy equivalence
$f:N'\hookrightarrow Z\to M$ consequently satisfies
\[
 \tau(f)=-b(x)+\overline{b(x)}=-b(x+x^*).                 \tag{A5.11.8}
\]
This uses the sign in Milnor's original Duality Theorem,
\cite[\S10, p.394, proof pp.395--397]{A5-Milnor1966}.
For clarity, one can check it on the two-term handle complex:
its dual has degrees $5\to4$, and torsion changes by the sign
$(-1)^6=+1$ and conjugate transpose.  No sign convention from
Siebenmann's Proposition 5.14 is needed.

Let $W\to N'$ be the infinite cyclic cover determined by
\[
 \pi_1(N')\xrightarrow{f_*}\pi_1(L)\times\mathbb Z
                       \longrightarrow\mathbb Z.
\]
The lift $\widetilde f:W\to L\times\mathbb R$ is a proper
homotopy equivalence: lift $f$, its compact-domain homotopy inverse,
and both homotopies.  It follows that $W$ has two ends, stable
fundamental group $C_{23}$ at each, and tame ends.  The geometric
infinite transfer in Example 11.5(ii), together with
Proposition 11.3(iii)--(iv), printed pp.123--125 of
\cite{A5-HughesRanicki1996}, identifies the end obstruction with
the Bass component of (A5.11.8), with the sign determined by the
choice of end.  In particular
\[
                  \sigma_+(W)=\pm(x+x^*)\ne0.           \tag{A5.11.9}
\]
Only nonvanishing is used.  Here is an additional algebraic check
fixing the sign for the convention in which positive deck powers
are truncated by $R[z]$.  The positive-end relative complex of
$\widetilde Z$ over $L\times\mathbb R$ has differential $A(z)$
in degrees $3\to2$, with cokernel $P$, and hence reduced Euler
class $x$.  The reverse complex has differential
$zA(z)^*=z(1-E^*)+E^*$ in degrees $5\to4$, with cokernel
$\ker(E^*)$ and reduced Euler class $-x^*$.  Shifting the chosen
lifts changes a truncated complex only by finite free complexes.
Additivity gives $\sigma_+(W)=x+x^*$ for this convention; at
the other end it is its negative.  This is also why the plus sign
in the proof on p.144 of \cite{A5-Siebenmann1972}, rather than
the inconsistent minus sign printed in Proposition 5.14 on p.143,
is compatible with this example.  The assertion itself is not
used as an input.

\paragraph{4. A torus product isomorphism and its bounded lift.}
The product cobordism $(Z\times S^1;M\times S^1,N'\times S^1)$
has zero torsion.  For example, filter the cellular chain complex
of the second circle by its one vertex and one edge: the two
relative complexes are base change of the handle complex of
$(Z,M)$ and its suspension, so their torsions cancel.
The smooth $s$-cobordism theorem, whose dimension hypothesis here
is amply satisfied, gives a diffeomorphism
\[
 h:N'\times S^1\xrightarrow{\cong}M\times S^1
                      =L\times S^1\times S^1             \tag{A5.11.10}
\]
homotopic to $f\times\operatorname{id}_{S^1}$.
We use \cite[\S10, $s$-Cobordism Theorem, p.393]{A5-Milnor1966}.
The homotopy ensures that $h$ preserves the specified homomorphism
of fundamental groups to $\mathbb Z^2$.

Let $\ell:W\to\mathbb R$ be a lift of the circle coordinate
$N'\xrightarrow f L\times S^1\to S^1$.
It is proper because the deck action is cocompact, and if $d$ is
the positive deck generator then $\ell(dw)=\ell(w)+1$.
Lifting (A5.11.10) gives a $\mathbb Z^2$-equivariant diffeomorphism
\[
 H:W\times\mathbb R\xrightarrow{\cong}
                   L\times\mathbb R\times\mathbb R,
 \qquad H(w,s)=(q,t,u).
\]
The functions $t-\ell(w)$ and $u-s$ are invariant under the
cocompact $\mathbb Z^2$ action.  Therefore some $B<\infty$ satisfies
\[
 |t-\ell(w)|\le B,\qquad |u-s|\le B
                         \quad\hbox{for every }(w,s).    \tag{A5.11.11}
\]
The same inequalities give the corresponding bounds for $H^{-1}$.
This uniform bound is the ingredient missing from an arbitrary
homeomorphism of stabilized open manifolds.

\paragraph{5. Extension over the distinguished line.}
Adjoin the positive end $v$ to $W$, obtaining the locally compact
metrizable space $C=W\cup\{v\}$.  Write
$cL=(L\times[0,\infty))/(L\times\{0\})$, with vertex $o$.
Set $\kappa(a)=\sqrt{1+a^2}$.  For $w\in W$ and $x\in\mathbb R$,
compute
\[
 (q,t,u)=H\bigl(w,\kappa(\ell(w))x\bigr),
 \qquad
 F(x,w)=\left(\frac{u}{\kappa(t)},[q,e^{-t}]\right),
 \qquad F(x,v)=(x,o).                                   \tag{A5.11.12}
\]
On the complements of the axes this is a homeomorphism, since
every coordinate change in the formula is a homeomorphism.
If $(x_i,w_i)\to(x,v)$ then $\ell(w_i)\to+\infty$.
The bounds (A5.11.11) imply $t_i\to+\infty$,
\[
 \frac{\kappa(\ell(w_i))}{\kappa(t_i)}\longrightarrow1,
 \qquad
 \frac{u_i}{\kappa(t_i)}
 =\frac{\kappa(\ell(w_i))}{\kappa(t_i)}x_i
       +\frac{u_i-\kappa(\ell(w_i))x_i}{\kappa(t_i)}
 \longrightarrow x.
\]
Since $L$ is compact, $[q_i,e^{-t_i}]\to o$ independently of
$q_i$.  This proves continuity at the axis.  Applying exactly the
same inequalities to $H^{-1}$ proves continuity of the inverse
there.  Thus (A5.11.12) proves the axis-fixing homeomorphism
(A5.11.1) directly.

\paragraph{6. Distinct transverse germs and maximality of the axis.}
If $(C,v)\cong(cL,o)$ as germs, the positive end of $W$ would
have a neighborhood homeomorphic to $L\times(0,\epsilon)$.
That end has zero finiteness obstruction, since $L$ is a finite
CW complex.  The end obstruction is a topological invariant and
is unchanged by passing to a smaller end neighborhood; this is
the necessity direction of the end theorem, discussed in
\cite[\S10, Proposition 10.2(i)--(ii), pp.110--111]{A5-HughesRanicki1996}.
It contradicts (A5.11.9).  This argument distinguishes germs,
not merely global spaces.

Every point of $C\setminus\{v\}=W$ is a manifold point, whereas
$v$ is not: a Euclidean germ at $v$ would likewise give a collared
end with zero obstruction.  The cone $cL$ also has exactly one
nonmanifold point, because
$H_1(L;\mathbb Z)=\mathbb Z/23$ gives a local homology group
in a degree other than its manifold dimension.  Both vertices
are immobile.  Indeed a positive-dimensional product axis through
an isolated nonmanifold point would contain other nonmanifold
points with the same germ, a contradiction.

For $X=\mathbb R\times cL$, its nonmanifold locus is exactly
$Y=\mathbb R\times\{o\}$: away from $Y$ it is a seven-manifold,
while on $Y$ the local homology satisfies
\[
 H_i(X,X\setminus\{(a,o)\};\mathbb Z)
             \simeq\widetilde H_{i-2}(L;\mathbb Z),
\]
so in particular $H_3=\mathbb Z/23$ there.  Any product axis
through a point of $Y$ lies in this same nonmanifold locus.
A manifold of dimension at least two cannot embed in a line;
for example an embedded small circle in such a manifold would
embed in the line, which is impossible after deleting a
non-extremal point.  Hence the mobility on $Y$ is at most one,
and the displayed product makes it exactly one.  Thus the two
normal factors in (A5.11.1) meet the precise maximal-axis and
immobility requirements of the bracket.

Finally $\mathbb R\times C$ is an ENR by (A5.11.1), because
$L$ is a compact smooth manifold.  Its retract $C$ is a locally
compact metrizable ANR, hence a finite-dimensional ENR; finite
dimensionality follows from its embedding as a factor slice in
the seven-dimensional product.  This verifies the usual local
regularity assumptions on the TOP spaces in the counterexample.

\paragraph{7. Exact finite certificate.}
The program \path{code/A5_top_normal_cancellation.py} checks
(A5.11.2)--(A5.11.6) using integer coefficient convolution in
$\mathbb Z[C_{23}]$ and an independent polynomial-remainder
implementation.  It also verifies the two quadratic ideal
identities by Hermite normal form.  Its complete output is in
\path{logs/A5_top_normal_cancellation.txt}; the full matrices and
coefficient vectors are in
\path{logs/A5_top_normal_cancellation.json}.  Selected output is
\begin{verbatim}
Gauss relation r^2-r+6=6*N: PASS.
E is an integral 2-by-2 idempotent of rank one: PASS.
Conjugate transpose E* is idempotent: PASS.
Independent polynomial-remainder verification: PASS.
Phi_23(g+1) is Eisenstein at 23: PASS.
p^3=(2-omega), both ideal lattices have index 8: PASS.
p*conjugate(p)=(2): PASS.
Image of x+x* is 2[J] != 0.
\end{verbatim}
The finite algebra certificate does not substitute for the
handle realization, the end theorem, or the continuity proof;
those inputs and their hypotheses were supplied above.
\end{proof}

\begin{saresult}{A5.012}{Proved}{The maximal triviality locus for proper maps with finite fibres}
Fix an o-minimal expansion of an arbitrary real closed field $K$.
Let $f:X\to Y$ be a continuous definable map of definable subsets
of finite-dimensional $K$-spaces, and assume it is \emph{definably
proper}: inverse images of definably compact sets are definably
compact. Suppose every fibre is finite. Then the maximal open
\emph{definable-manifold} locus $U\subset Y$ over which $f$
is definably locally trivial is definable, and $Y\setminus U$
is nowhere dense. Manifolds here have no boundary; no triangulation
of the map is required.
\end{saresult}
\begin{proof}
Let $Y_{\rm man}$ be the open locus where $Y$ is locally a
definable manifold without boundary. This set is definable by the
finite pointed-germ argument in A5.005, which uses only finite
definable triangulation and applies over $K$ as well. Replace $Y$
by $Y_{\rm man}$ and $X$ by its inverse image; definable properness
is preserved. Let
$L\subset X$ be the locus where $f$ is a local homeomorphism
onto an open subset of $Y$. We first show directly that $L$
is definable, without quantifying over homeomorphisms.
For $x\in X$, write $V_r(x)=X\cap B(x,r)$. Membership in $L$
is equivalent to the existence of $r>0$ such that $f$ is injective
on $V_r(x)$ and the following condition holds:
\[
\begin{gathered}
 \forall z\in V_r(x)\quad\forall\delta>0\quad
 \exists\varepsilon>0\quad\forall y\in Y,\\
 \|y-f(z)\|<\varepsilon\quad\Longrightarrow\\
 \exists w\in V_r(x):\quad\|w-z\|<\delta
       \quad\text{and}\quad f(w)=y.
\end{gathered}                                      \tag{A5.12.1}
\]
Indeed (A5.12.1) says exactly that the restriction of $f$ to
$V_r(x)$ is open into $Y$. Together with injectivity and the
given continuity, this makes it a homeomorphism onto an open
subset. Conversely, a local homeomorphism has such a ball
inside its domain of definition. All quantifiers in (A5.12.1)
range over $K$-coordinates and radii in definable sets; hence
it defines a definable set. Its definition also shows that $L$
is open in $X$.

We claim that the desired locus, in the original target, is
\[
 U=Y_{\rm man}\setminus f\bigl(f^{-1}(Y_{\rm man})\setminus L\bigr).
                                                               \tag{A5.12.2}
\]
The right side is definable and open. Here is the closed-image
argument over $K$, where ordinary topological compactness cannot
be substituted for definable compactness. Let $C$ be a closed
definable subset of the source and $y\in\overline{f(C)}^{\,Y}$.
By definable curve selection choose a continuous definable curve
$\alpha:[0,\varepsilon]\to Y$ with $\alpha(0)=y$ and
$\alpha(t)\in f(C)$ for $t>0$. The image
$Q=\alpha([0,\varepsilon])$ is definably compact.
Definable choice supplies a lift $\beta:(0,\varepsilon]\to C$
with $f\circ\beta=\alpha$, continuous after shortening the
interval. Its image lies in the definably compact set $f^{-1}(Q)$.
Coordinatewise monotonicity gives a finite endpoint limit
$x=\lim_{t\to0^+}\beta(t)$ in that set. Since $C$ is closed
in the source, $x\in C$; continuity gives $f(x)=y$.
Thus $f(C)$ is closed. Applying this to the displayed bad set
proves openness of (A5.12.2).

If $f$ is locally trivial near $y$, its finite fibre is discrete,
so every point of $f^{-1}(y)$ belongs to $L$.
Conversely let $y$ lie in (A5.12.2), and write
$f^{-1}(y)=\{x_1,\ldots,x_k\}$. Choose pairwise disjoint definable
open neighbourhoods $V_i$ of $x_i$ such that
$f_i=f|_{V_i}:V_i\to W_i$ is a homeomorphism onto a definable
open neighbourhood of $y$. By the definable closed-image argument the image of
$X\setminus\bigcup_iV_i$ is closed and misses $y$.
Choose a definable open neighbourhood
\[
 W\subset\bigcap_iW_i\setminus
                  f\bigl(X\setminus\bigcup_iV_i\bigr)
\]
of $y$. The typed map
\[
 H:f^{-1}(y)\times W\longrightarrow f^{-1}(W),\qquad
                  H(x_i,z)=f_i^{-1}(z)
\]
is a definable homeomorphism and $f\circ H=\operatorname{pr}_W$.
For an empty fibre use the open complement of the closed image
of $f$ and the empty-product trivialization. Thus (A5.12.2)
is precisely the maximal locus. Its complement is nowhere dense
by A5.008. This argument settles all fibre-dimension-zero cases;
the remaining general gate concerns positive-dimensional fibres.
\end{proof}


\begin{saresult}{A5.013}{Proved under stated hypotheses}{The maximal triviality locus for targets of dimension at most two}
Let $M$ be an o-minimal expansion of a real closed field and let
$f:X\to Y$ be a continuous definably proper definable map.  Suppose
$\dim Y\leq2$.  Then the largest open subset of $Y$ which is a
definable manifold without boundary and over which $f$ is definably
locally trivial is definable.  Thus the remaining maximal-locus
question requires a target of dimension at least three and
positive-dimensional fibres.
\end{saresult}
\begin{proof}
All neighbourhoods and maps used below are definable; parameters
are allowed.  ``Compact'' in this proof means closed and bounded,
or equivalently definably compact, and never asserts ordinary
topological compactness over a non-Archimedean field.

\medskip\noindent\textbf{1. A finite-format consequence of logical compactness.}
We first give the precise uniformity principle used twice below.
Fix the structure $M$, parameters defining a map $f$, and a
definable set $A$ of target points.  A \emph{certificate format}
is a finite list of formulas for sets and graphs, together with
fixed finite combinatorial data.  Its free tuple $b$ supplies all
geometric parameters.  A predicate $C(a,b)$ is first-order if it
requires the indicated graphs to be continuous maps or
homeomorphisms on the indicated domains, requires specified
diagram identities, and requires the domains to be open, closed,
bounded, or to contain specified points or intervals.  For
example, continuity of a graph $\Gamma_h$ at $x$ is the formula
\[
 \forall\epsilon>0\ \exists\delta>0\ \forall x',z,z'
 \bigl((x,z),(x',z')\in\Gamma_h\ \wedge\ |x-x'|<\delta
       \Longrightarrow |z-z'|<\epsilon\bigr).
\]
Single-valuedness, surjectivity, injectivity and inverse continuity
are likewise first-order.  A fixed finite polyhedron can be
represented by its vertex coordinates and a fixed list of faces;
the required incidence and affine-independence conditions are
finite first-order conditions.  A finite piecewise affine map,
possibly discontinuous on the boundaries of its pieces, has the
same kind of description.

Suppose that in every elementary extension $M'\succeq M$, every
$a\in A(M')$ has a certificate of one of a set of formats
$\{C_i(a,b_i):i\in I\}$.  Then finitely many formats suffice.
Indeed, if no finite list sufficed, the elementary diagram of $M$
together with a new constant $a\in A$ and the formulas
$\neg\exists b_i\,C_i(a,b_i)$ for all $i\in I$ would be finitely
satisfiable in $M$.  The compactness theorem would give an
elementary extension containing an $a$ with no certificate,
contrary to the hypothesis.  The compactness theorem used here is
Theorem~2.4, p.17 of Marker's original Orsay notes
\cite{A5-Marker2010}; this paragraph supplies the entire
application, including the elementary-diagram argument.

Every elementary extension here is still o-minimal.  To see why
this use is legitimate, apply cell decomposition to each fixed
one-variable formula with parameters in $M$.  It gives a uniform
finite bound on the number of interval and point components.
The assertion that every parameter fibre is a union of at most
that many intervals and points is a first-order sentence, using
their endpoints.  It holds in the extension, for each formula.
All closed-boundedness and continuity hypotheses used below also
transfer by their first-order formulations.

Here is the particular consequence concerning local PL models.
Assume that in every elementary extension and at every $a\in A$
there is a local isomorphism of the full map over a target
neighbourhood to a finite simplicial map
$p:|K|\to|L|$.  More explicitly, the certificate consists of
open sets $V\subset Y$, $W\subset|L|$, homeomorphisms
\[
 \beta:W\xrightarrow{\sim}V,
 \qquad
 \alpha:p^{-1}(W)\xrightarrow{\sim}f^{-1}(V),
 \qquad f\alpha=\beta p,
 \qquad a\in V.
 \tag{A5.13.1}
\]
The finite complexes and their vertex map are part of the format;
use their canonical rational realizations, so there are no
varying combinatorial types within one format.  All requirements
in (A5.13.1) are first-order.  The preceding argument therefore
gives finitely many such chart formats covering $A$.

For each of their fixed simplicial maps $p_i$, let $G_i\subset|L_i|$
be its manifold-and-local-triviality locus in the original structure
$M$.  No invariance of that unknown goodness property under
elementary extensions is asserted or needed.  By A5.008,
$G_i$ is a union of finitely many open simplices, and hence is
semilinear.  No decision procedure for choosing these finitely
many simplices is required.  In any certificate of the $i$th
format, local triviality and manifoldhood at $a$ are equivalent
to membership of $\beta^{-1}(a)$ in this same $G_i$.  The formula
\[
 a\in A\ \wedge\
 \bigvee_i\exists b_i\,\exists w\,
 \bigl(C_i(a,b_i)\ \wedge\ (w,a)\in\Gamma_{\beta_i,b_i}
                  \ \wedge\ w\in G_i\bigr)
 \tag{A5.13.2}
\]
thus defines the desired locus on $A$.  This argument does not
assume that the unknown good locus has uniform chart complexity:
it only uses universal existence of PL models on the already
definable set $A$.

\medskip\noindent\textbf{2. A semilinear agreement gives an actual local PL model.}
Let $P$ be a compact finite polyhedron, let $F:P\to Z\subset M^m$
be continuous definable, and let $g:P\to M^m$ be a finite
piecewise affine map, not assumed continuous everywhere.  Suppose
$V\subset Z$ is open and
\[
             F(x)=g(x)\qquad(x\in F^{-1}(V)).
 \tag{A5.13.3}
\]
At every $y\in V$ there is a local finite PL model of the full
map $F$, provided the target has a PL chart near $y$.

Write $O=F^{-1}(V)$ and $E=F^{-1}(y)$.  The equality
$E=O\cap g^{-1}(y)$ shows that $E$ is open in the semilinear set
$g^{-1}(y)$; it is also closed there because $E$ is closed in $P$.
It is a definable clopen subset of that semilinear set.  Each
semilinear definably connected component is semilinearly path connected,
so a definable clopen subset is a union of such components.
There are finitely many components.  Consequently $E$ is a
compact semilinear subpolyhedron of $P$.

If $E$ is empty, compactness of $P$ makes $F(P)$ closed and the
empty map over a small neighbourhood is a finite PL model.
Otherwise triangulate $P$ compatibly with $E$, making the latter
a full subcomplex after barycentric subdivision.  The PL function
$\rho:P\to[0,1]$ which is zero on vertices of $E$ and one on the
other vertices has zero set exactly $E$.  On the compact set
$P\setminus O$ its positive values have a positive minimum when
that set is nonempty.  Choose $c>0$ below that minimum (and any
$c>0$ if $P\setminus O$ is empty), and put
\[
                     Q=\{x\in P:\rho(x)\leq c\}.
\]
Then $Q$ is a compact polyhedral neighbourhood of $E$, contained
in $O$.  On $Q$, $g=F$ is continuous and piecewise affine, hence
is a finite PL map.  The compact set $P\setminus\operatorname{Int}_P Q$
has closed image under $F$, disjoint from $y$.  A small target
neighbourhood $V_y\subset V$ avoids this image, so
$F^{-1}(V_y)\subset\operatorname{Int}_P Q$.  Thus
$F:F^{-1}(V_y)\to V_y$ is exactly the restriction of the finite
PL map $g|_Q$, with its entire fibre.  In the planar target chart,
take a small closed box $B$ about $y$ contained in that chart.
The union $S=g(Q)\cup B$ is a compact target polyhedron; shrink
$V_y$ further into the interior of $B$.  Triangulate the finite
PL map $g|_Q:Q\to S$ (which need not be surjective), and pass
to the canonical realizations of these finite complexes to obtain
(A5.13.1).  The same construction works in any finite-dimensional
PL chart.  All arguments in this step work in every elementary
extension.

\medskip\noindent\textbf{3. A punctured edge admits semilinear agreement.}
The input from Shiota is the unnumbered statement on printed p.33
of \cite{A5-Shiota}, proved on pp.33--34: if $K$ is a finite
simplicial complex and
$H=(H_1,H_2):|K|\to M^2$ is continuous definable with both
coordinates nonnegative, there are a definable isotopy $\zeta_t$
of $|K|$ preserving each simplex, a number $\epsilon>0$, and a
definable relative neighbourhood $N$ of
$(0,\epsilon]\times\{0\}$ in
$(0,\epsilon]\times[0,\epsilon]$, such that
$H\zeta_1$ on $(H\zeta_1)^{-1}(N)$ extends to a PL map
$L:|K|\to M^2$.  Shrinking $N$ makes it relatively open and
restricts to $0<s<\epsilon$, $0\leq t<\epsilon$.

Consider a continuous map from a compact definable set to a
two-dimensional polyhedral target, an interior point $a$ of a
target edge, and the case in which that edge has exactly two
incident triangles.  Affinely flatten the two triangles onto
opposite sides of an axis; in a neighbourhood of $a$ use
coordinates $(s,t)$ with the edge given by $t=0$.  For a sufficiently
small $b>0$, the closed right rectangle
$D=[a,a+b]\times[-b,b]$ lies in this chart.  Its full preimage
$X_D$ is compact.  Triangulate $X_D$ compatibly with the closed
subsets given by $t\circ f\geq0$ and $t\circ f\leq0$:
\[
 h_0:P=|K|\xrightarrow{\sim}X_D,
 \qquad P_+,P_-\subset P,
 \qquad P_0=P_+\cap P_- .
\]
These three subsets are subcomplexes.  Apply Shiota's statement to
\[
 H=(s\circ f\circ h_0-a,\ |t\circ f\circ h_0|):P\longrightarrow M^2.
\]
Because $\zeta_1$ preserves the subcomplexes, the homeomorphism
$h=h_0\zeta_1:P\to X_D$ preserves their interpretation as the two
sides.  Define a finite piecewise affine map $g:P\to M^2$ by
\[
 g(x)=
 \begin{cases}
 (a+L_1(x), L_2(x)),&x\in P_+\setminus P_0,\\
 (a+L_1(x),-L_2(x)),&x\in P_-\setminus P_0,\\
 (a+L_1(x),0),&x\in P_0.
 \end{cases}
 \tag{A5.13.4}
\]
It may fail to be continuous outside the neighbourhood in question;
Step~2 was formulated to allow exactly this possibility.
The signed inverse image of $N$ under $(s,t)\mapsto(s-a,|t|)$,
restricted to the interior of $D$, is an open set $V$ containing
$(a,a+\delta)\times\{0\}$ for some $\delta>0$.  On
$(f h)^{-1}(V)$, equation (A5.13.4) agrees with the signed map
$f h$.  Indeed, the unsigned coordinates agree with $L$ there,
the side is preserved, and $L_2=0$ on $P_0$ there.  Step~2 gives
a local finite PL model at every point of this deleted right
edge interval.  The left interval is obtained by using $a-s$.

Retain the stronger data, not merely the conclusion: a rectangle
$D$, a polyhedron $P$, a homeomorphism $h:P\to f^{-1}(D)$,
a finite piecewise affine map $g$, a target open set
$V\subset\operatorname{Int}D$, and a number $\delta>0$, such that
$V$ contains the specified deleted interval and $fh=g$ on
$(fh)^{-1}(V)$.  For each fixed choice of formulas and finite
piece data these are first-order requirements, of the kind
described in Step~1.

Fix one open target edge $e$.  At every $a\in e$, in every
elementary extension, the preceding construction supplies both
left and right certificates.  Applying Step~1 to pairs of such
formats yields finitely many formats covering all $a\in e$.
Let $D_e\subset e$ be the union of all deleted intervals certified
by all parameter instances of these finitely many formats.
This is a definable set: membership is an existential formula
with a finite disjunction.  For every $a\in e$ there is a
$\delta_a>0$ for which
\[
       (a-\delta_a,a+\delta_a)\cap e\setminus\{a\}\subset D_e.
 \tag{A5.13.5}
\]
The definable complement $e\setminus D_e$ has no interval, and
therefore is finite.  Every point of $D_e$ has a local finite PL
model by Step~2, in every elementary extension, because the
certificates asserting (A5.13.3) transfer.  Apply the local-model
consequence (A5.13.2) of Step~1 to the already definable set
$A=D_e$.  It proves definability of the maximal triviality locus
on $D_e$.  Its value on the remaining finite set $e\setminus D_e$
is also definable with parameters.  The desired locus on the
entire open edge is consequently definable.

\medskip\noindent\textbf{4. Reduction to the compact planar situation.}
Let $Y_{\rm man}$ denote the definable manifold-without-boundary
locus of $Y$.  Its definability was proved in A5.005.  The
components of local dimension zero or one are settled by A5.008.
Restrict to the union $T$ of its components of dimension two and
to $f^{-1}(T)\to T$.  This map remains definably proper.  The
set $T$ is locally compact and locally closed in its ambient
space.  One direct verification of the latter fact is to take
at each point a closed small ball inside a definable manifold
chart; its image is compact and is an ambient closed
neighbourhood in $T$.  Thus $T$ is open in its ambient closure.

Apply coordinatewise semialgebraic bounded homeomorphisms to
source and target, and let $\overline\Gamma$ be the closure of
the resulting graph inside a closed bounded box.  Put
$\overline T=\operatorname{cl}T$ in the bounded target coordinates.
Projection gives a continuous map of compact definable sets
\[
          \overline f:\overline\Gamma\longrightarrow\overline T.
\]
There are no added graph points over $T$.  In fact, at any
$y\in T$ take a compact target neighbourhood $C\subset T$ whose
interior contains $y$.  Properness makes its original preimage
compact; its transformed graph is compact and closed.  Any
point of $\overline\Gamma$ over $y$ is in its closure, since the
target coordinate of a convergent definable curve eventually
lies in the interior of $C$.  It therefore belongs to the
original graph.  This also proves
$\overline f^{-1}(T)\cong f^{-1}(T)$ over $T$.

Take a Hardt partition of $\overline f$, and a finite definable
triangulation of $\overline T$ compatible with that partition
and with $T$.  Its dimension is at most two.  Transfer the map
to this polyhedral target.  On every open two-simplex contained
in $T$, the full map is definably trivial by the chosen Hardt
partition, so all its points belong to the maximal open locus.
An open one-simplex contained in the two-manifold $T$ has exactly
two incident two-simplexes; its usual finite-polyhedral
neighbourhood is therefore the planar chart used in Step~3.
That step proves definability of the locus on every such edge.
There are finitely many remaining vertices.  Any subset of them
is definable with parameters.  The union over the finitely many
simplexes is the full maximal locus in $T$, and hence is
definable.  Joining the already settled components of dimensions
zero and one proves the theorem.
\end{proof}


% Independently audited by root, PASS, 2026-11-19 16:27 UTC.
\begin{saresult}{A5.014}{Proved under stated hypotheses}{The maximal triviality locus for compact graph fibres}
Let $M$ be an o-minimal expansion of a real closed field, and let
$f:X\to Y$ be a continuous definably proper definable map whose
fibres have dimension at most one.  Then its maximal open
definable-manifold-without-boundary local-triviality locus is
definable, with no restriction on $\dim Y$.
For an arbitrary continuous definably proper map, the part of its
maximal locus over which the fibres have dimension at most one
is also definable.  Combined with A5.013, the remaining question
requires fibre dimension at least two and target dimension at
least three.
\end{saresult}
\begin{proof}
We give a definable construction over the given field; ordinary
compactness of metric spaces and integration of functions are
not used.  Replace $Y$ by its definable manifold locus, which is
open and definable by A5.005.  Embed $X$ in $M^n$, and write
$F_y=f^{-1}(y)$.  Each fibre is a compact finite graph, allowing
isolated points, empty graphs, and disconnected graphs.

\medskip\noindent\textbf{1. Definable arc length and a uniform replacement estimate.}
For a definable embedded arc $\gamma:[a,b]\to M^n$, define its
length by
\[
 \ell(\gamma)=\sum_{j=1}^n\operatorname{Var}_{[a,b]}(\gamma_j).
 \tag{A5.14.1}
\]
Every coordinate has finitely many monotonicity intervals.  If
their number is at most $N$, its variation is the maximum of
\[
 \sum_{i=1}^{N+1}|\gamma_j(t_i)-\gamma_j(t_{i-1})|,
 \qquad a=t_0\leq t_1\leq\cdots\leq t_{N+1}=b.
\]
Repeated partition points are permitted.  This is a maximum on
a compact definable simplex, and is first-order definable.
Length is additive under concatenation, independent of an
increasing reparameterization, and strictly positive on every
nonconstant subarc.  Cumulative length is a continuous strictly
increasing definable function, and hence parameterizes an arc
by an interval $[0,\ell]$.  Ambient distance is bounded by this
length.  A definable embedded circle similarly has a definable
total length $L>0$ and shorter-arc metric $d$; it is the metric
circle of circumference $L$ in these coordinates.

There is a single finite bound $N$ adequate for all the fibres
in the present definable family.  Indeed, Hardt triviality
\cite[Theorem 5.22, pp.62--63]{A5-Coste}, followed by
triangulation of its finitely many compact representative
fibres, gives finitely many definable formula families
$\gamma_e(y,t)$ parameterizing their edges.  Monotonicity and
cell decomposition, with $t$ the last coordinate, give a uniform
bound for every ambient coordinate of every such edge
\cite[Theorems 2.1, 2.9--2.10, pp.15, 19]{A5-Coste}.  Coalescing
bivalent subdivisions multiplies this bound by at most the fixed
number of edges.  Cutting a circle at an arbitrary point adds
at most one break.  Enlarge $N$ once to cover all these cases.

An \emph{intrinsic vertex} of a graph has local valence different
from two; isolated points have valence zero.  The components
left after removing intrinsic vertices are its intrinsic edges.
A component without vertices is a circle.  A homeomorphism
permutes this finite intrinsic graph, is monotone on each
oriented edge, and has one of the two orientations on a circle.
Replace an edge homeomorphism by the map linear in normalized
length coordinates.  On a circle prescribe the image of one
base point and make the same replacement around the circle.
Call these maps \emph{canonical}.

Suppose $h:F_y\to F_z$ is a definable homeomorphism satisfying
$\|h(x)-x\|<\epsilon$ for all $x$.  On one source edge, transport
the target coordinate functions along $h$.  Each coordinate of
both curves has at most $N$ monotonicity intervals.  Refining
their two partitions gives at most $2N+2$ intervals.  On each
interval the absolute difference between the two coordinate
variations is at most $2\epsilon$, by comparing the two
endpoints.  Thus the source prefix length $u$ and target prefix
length $v$ of $h(x)$ satisfy
\[
       |u-v|\leq C\epsilon,
       \qquad |\ell_y-\ell_z|\leq C\epsilon,
       \qquad C=4n(N+1).
 \tag{A5.14.2}
\]
The canonical replacement sends $u$ to
$v'=u\ell_z/\ell_y$.  Since $0\leq u\leq\ell_y$,
$|v'-u|\leq C\epsilon$, so $|v'-v|\leq2C\epsilon$.
Ambient distance is bounded by length.  Consequently its
displacement is less than
\[
                         (1+2C)\epsilon.                 \tag{A5.14.3}
\]
The argument applies to a cut circle with base point image
$h(p)$.  For a loop based at an intrinsic vertex, use the
interval parametrization whose two endpoints are identified
at that vertex; the same variation and prefix-length estimates
apply to this cut parametrization.  At isolated vertices
the original displacement estimate suffices.

All canonical maps occur in finitely many first-order formula
formats.  The finite data specify the Hardt pieces of $y,z$,
an isomorphism of their intrinsic graphs, and the orientations
of the edges.  Only the images of the chosen circle base points
are additional geometric parameters.  Formula (A5.14.1) and
the unique inverse of cumulative length define the map on each
edge; the definitions agree at the vertices.  For circles use
length modulo their circumference.  Thus the graph, its
homeomorphism property, and its displacement bound are all
expressible by finitely many formulas.  This conclusion uses
the one-dimensional monotonicity estimate, not an assumed
uniform bound for arbitrary graph homeomorphisms.

\medskip\noindent\textbf{2. Complete regularity is definable.}
Call $y$ completely regular if for every $\epsilon>0$ there is
$\delta>0$ such that every $z\in Y$ with $\|z-y\|<\delta$
has a definable homeomorphism $h:F_y\to F_z$ of displacement
less than $\epsilon$.  Step~1 shows that the same condition is
obtained by requiring $h$ to be canonical: use
$\epsilon/(1+2C)$ before applying (A5.14.3).  Its precise
first-order form is therefore
\[
 \begin{gathered}
 \forall\epsilon>0\ \exists\delta>0\ \forall z\in Y,\quad
 \|z-y\|<\delta\ \Longrightarrow\\
 \bigvee_i\exists b_i\,
 \left(C_i(y,z,b_i)\ \wedge\
 \forall(x,x')\in\Gamma_{i,y,z,b_i},\ \|x'-x\|<\epsilon\right),
 \end{gathered}                                           \tag{A5.14.4}
\]
where $C_i$ specifies a canonical homeomorphism and
$\Gamma_i$ its graph.  Let $R_f\subset Y$ be the definable set
given by (A5.14.4).

If $f$ is locally trivial at $y$, make its product chart the
identity on $F_y$.  Uniform continuity of this chart on the
product of $F_y$ with a compact small target neighbourhood
gives (A5.14.4).  The same applies at every point of that
neighbourhood.  Hence the maximal locus is contained in
$\operatorname{Int}_Y R_f$.  It remains to prove the converse.

\medskip\noindent\textbf{3. The metric circles vary continuously on $R_f$.}
Fix a point $y_0\in\operatorname{Int}R_f$ and write $F_0=F_{y_0}$.
The empty-fibre case is immediate from (A5.14.4), which then
forces all nearby fibres to be empty.  Otherwise shrink to an
open neighbourhood $V$ contained in $\operatorname{Int}R_f$
such that every fibre is canonically homeomorphic to $F_0$
by a map of a fixed sufficiently small displacement.  Distinct
components of $F_0$ have positive separation.  Thus its circle
components can be labelled throughout $V$ by disjoint fixed
ambient neighbourhoods.  Their unions in $f^{-1}(V)$ are both
open and closed relative to $f^{-1}(V)$, since the displacement
bound leaves a positive gap between those neighbourhoods and
the other components.

For one such circle family $C_y$, denote its length and metric
by $L_y,d_y$.  They are definable by Step~1.  They are also
continuous, with $d_y$ jointly continuous on the fibre-pair
family.  In fact, at any $y\in V$, complete regularity provides
canonical comparisons $k:C_y\to C_z$ with arbitrarily small
displacement as $z\to y$.  Equation (A5.14.2) shows
$L_z\to L_y$.  Since the comparison is linear in normalized
length,
\[
 d_z(k(x),k(x'))=\frac{L_z}{L_y}d_y(x,x').                 \tag{A5.14.5}
\]
For arbitrary $x_z,x'_z$ approaching $x,x'$, their inverse
images under $k$ approach $x,x'$ in the fixed compact circle
$C_y$, because $k^{-1}$ has the same displacement bound as
$k$.  The metric $d_y$ has the ordinary definable topology
on this fixed circle and is uniformly continuous.  Equation
(A5.14.5) proves the assertion by epsilon-delta estimates.
This reasoning is valid in the given ordered field; no
sequential compactness is being substituted for definable
compactness.

\medskip\noindent\textbf{4. Continuous definable anchors for the circle components.}
Choose a regular $C^1$ point $p\in C_0$; such points exist by
definable $C^1$ cell decomposition
\cite[Theorem 6.6, p.71]{A5-Coste}.  There is a constant
$A>0$ such that
\[
                     d_0(p,z)\leq A\|z-p\|\qquad(z\in C_0).
 \tag{A5.14.6}
\]
Locally at $p$, choose one coordinate with nonzero tangent
derivative and use it as the arc parameter.  The other
coordinate derivatives are bounded on a smaller compact arc,
so its $L^1$ length is bounded by a constant times that
coordinate difference.  Away from the arc, compactness gives
a positive lower bound for $\|z-p\|$, whereas $d_0\leq L_0/2$.
These two estimates prove (A5.14.6).

For a constant $M>A^2$ to be enlarged below, define
\[
 E_y(x)=\max_{z\in C_y}
             \left(d_y(x,z)^2-M\|z-p\|^2\right).
 \tag{A5.14.7}
\]
The maximum exists and is definable.  By (A5.14.6),
$E_0(p)=0$.  The term $z=p$ gives
$E_0(x)\geq d_0(x,p)^2$, so $p$ is its unique global
minimizer.  Under a small canonical comparison
$k:C_0\to C_y$, (A5.14.5), $L_y\to L_0$, and
$\sup\|k(z)-z\|\to0$ show that
$E_y(k(x))\to E_0(x)$ uniformly in $x$.  This follows directly
by bounding the difference inside the maximum in (A5.14.7);
all lengths and ambient distances are uniformly bounded on
the compact reference circle and its small neighbourhood.

Choose concentric short closed intrinsic arcs $J_0\subset
\operatorname{Int}I_0$ about $p$, with the radius of $I_0$
less than $L_0/8$, and $J_0$ strictly smaller.  Points of
$C_0\setminus\operatorname{Int}I_0$ have ambient distance
at least $m>0$ from $p$.  Enlarge $M$ so that, for all nearby
circles, a point at distance at least $m/2$ from $p$ has
\[
            d_y(x,z)^2-M\|z-p\|^2<-1
\]
for every $x\in C_y$; use a fixed upper bound for $L_y/2$.
The point $k(p)$, on the other hand, contributes a value
greater than $-1/2$ when $y$ is sufficiently close to $y_0$.
Thus every maximizer defining (A5.14.7) lies in $I_y=k(I_0)$.
The strict minimum and positive gap of $E_0$ outside
$\operatorname{Int}J_0$, together with the preceding uniform
estimate, show that every global minimizer of $E_y$ lies in
$\operatorname{Int}J_y$, where $J_y=k(J_0)$.

For $x\in J_y$ and $z\in I_y$ the shorter arc has no
wrap-around.  In an intrinsic length coordinate $u$ on
$I_y$, equation (A5.14.7) becomes
\[
 E_y(u)=u^2+\max_{z\in I_y}
       \left(-2u\,u(z)+u(z)^2-M\|z-p\|^2\right).
 \tag{A5.14.8}
\]
The second term is convex as a maximum of affine functions.
More explicitly, for $0<t<1$ its values satisfy
\[
 E_y((1-t)u+tv)\leq(1-t)E_y(u)+tE_y(v)-t(1-t)(u-v)^2.
\]
Consequently $E_y$ is strictly convex on $J_y$.  Since all
global minimizers lie there, there is exactly one; denote it
by $s(y)$.  The equation of a global minimizer is a first-order
condition, so $s$ is definable.

It is continuous throughout a smaller $V$.  At any $y\in V$,
small canonical comparisons and (A5.14.5) make the energies
uniformly close on the fixed reference circle $C_y$.  That
circle has a unique minimizer of its continuous energy.
Outside any chosen neighbourhood of that minimizer, compactness
gives a positive gap above the minimum.  Uniform closeness
therefore puts the minimizer for every sufficiently close
$z$ in the corresponding neighbourhood.  Small displacement
of the comparison makes this an ambient neighbourhood as well.
This proves continuity by epsilon-delta estimates.  Repeating
the construction for the finitely many circle components
supplies continuous definable anchors $s_j(y)$, with
$s_j(y_0)=p_j$.

\medskip\noindent\textbf{5. Unique small canonical maps produce a bundle.}
Mark $F_0$ at the chosen point $p_j$ on each circle component.
Its canonical automorphisms fixing these marked points form
a finite group $G$: the intrinsic finite graph has only
finitely many permutations and orientations; on a marked
circle there are just the identity and reflection.  Choose
$\eta>0$ so that
\[
 2\eta<\min_{a\in G,\ a\ne1}\max_{x\in F_0}\|a(x)-x\|,
 \tag{A5.14.9}
\]
and choose any $\eta>0$ if $G$ is trivial.  The positive
minimum exists because the set is finite and every nonidentity
map moves some point.

After shrinking $V$, there exists for every $y\in V$ a
canonical homeomorphism $h_y:F_0\to F_y$ with
$h_y(p_j)=s_j(y)$ and displacement less than $\eta$.
To see this, begin with an arbitrarily small canonical map
given by complete regularity at $y_0$.  On each target circle,
rotate its length coordinate to send the image of $p_j$ to
$s_j(y)$.  The required intrinsic rotation tends to zero:
transport back to the fixed circle by (A5.14.5), and use
continuity of its metric and of $s_j$.  Ambient motion is
at most the intrinsic rotation.  The adjusted canonical map
therefore has displacement less than $\eta$.

It is unique.  If $h,h'$ were two such maps, then
$a=h^{-1}h'\in G$, and for every $x\in F_0$,
\[
 \|a(x)-x\|\leq\|a(x)-h(a(x))\|+\|h'(x)-x\|<2\eta.
\]
Equation (A5.14.9) forces $a=1$.  The finite formula formats
from Step~1, together with the anchor conditions and the
strict displacement bound, consequently define the graph of
a single family $h_y$.

This family is continuous.  Fix $y\in V$.  Complete regularity
at $y$ gives canonical comparisons $k:F_y\to F_z$ of
arbitrarily small displacement for $z$ near $y$.  For small
enough displacement the labelled circle components match.
Rotate them as above to make $k(s_j(y))=s_j(z)$; by continuity
of the anchors and (A5.14.5), this changes the displacement
by a quantity tending to zero.  Call the result $k'$.  The
composite $k'h_y$ is canonical, respects all anchors, and has
displacement below $\eta$ for $z$ sufficiently close to $y$,
since $h_y$ has a strict margin below $\eta$.  Uniqueness gives
\[
                             h_z=k'h_y.
\]
Hence $h_z\to h_y$ uniformly in the ambient metric.  This
implies joint continuity of
\[
 H:F_0\times V\longrightarrow f^{-1}(V),\qquad
 H(x,y)=h_y(x),\qquad fH=\operatorname{pr}_V.
 \tag{A5.14.10}
\]
The map is a definable continuous bijection.  It is definably
proper: the preimage of a compact set $A\subset f^{-1}(V)$
is a closed subset of the compact set $F_0\times f(A)$.
The definable closed-image argument from A5.012 makes its
inverse continuous.  Thus (A5.14.10) is a local definable
bundle chart.

We have proved that the maximal locus is exactly
$\operatorname{Int}_Y R_f$, which is definable.  Finally,
for an arbitrary proper map let
$Z=\{y:\dim F_y\leq1\}$.  A locally trivial value with such
a fibre lies in $\operatorname{Int}_Y Z$.  Restrict the map
to that definable open set and apply the result just proved.
This gives precisely the claimed low-dimensional-fibre part
of the maximal locus and proves the final reduction.
\end{proof}
